\documentclass[11pt]{article}

% ============================================================================
%  Elementary reductions for the two-powers squarefree variant of
%  Erdos Problem #11.
%
%  Self-contained research note. Compiles with pdflatex/xelatex/tectonic
%  (Overleaf or any standard TeX distribution); uses only widely available
%  packages.
%
%  HONESTY POLICY: only results that are fully proved here, or proved
%  elsewhere and cited exactly (Kalinin 2026), are stated as theorems.
%  Everything empirical-only, conditional, or still missing a rigorous
%  closure is confined to Section 7 and explicitly flagged as open.
% ============================================================================

\usepackage[T1]{fontenc}
\usepackage[utf8]{inputenc}
\usepackage{amsmath,amssymb,amsthm}
\usepackage{booktabs}
\usepackage[margin=1in]{geometry}
\usepackage{hyperref}
\hypersetup{colorlinks=true,linkcolor=blue,urlcolor=blue,citecolor=blue}

\theoremstyle{plain}
\newtheorem{theorem}{Theorem}[section]
\newtheorem{lemma}[theorem]{Lemma}
\newtheorem{proposition}[theorem]{Proposition}
\newtheorem{corollary}[theorem]{Corollary}
\theoremstyle{definition}
\newtheorem{definition}[theorem]{Definition}
\newtheorem{example}[theorem]{Example}
\theoremstyle{remark}
\newtheorem{remark}[theorem]{Remark}

\newcommand{\Z}{\mathbb{Z}}
\newcommand{\F}{\mathbb{F}}
\newcommand{\ord}{\operatorname{ord}}
\newcommand{\Np}{N_p(n)}
\newcommand{\dpp}{d_p}
\newcommand{\ep}{e_p}
\renewcommand{\proofname}{Proof}

\title{\bf Elementary reductions for the two-powers squarefree\\
variant of Erd\H{o}s Problem~\#11}
\author{Research note}
\date{23 June 2026 (revised 3 July 2026)}

\begin{document}
\maketitle

\begin{abstract}
\noindent Erd\H{o}s Problem~\#11 asks whether every sufficiently large integer
$n$ can be written as $n=k+2^m$ with $k$ squarefree. We study a relaxation,
the \emph{two-powers variant}: every odd $n>1$ satisfies
$n=k+2^l+2^m$ for some squarefree $k$ and $l,m\ge0$. This variant is not one of
Erd\H{o}s's numbered problems; it appears as the statement
\texttt{erdos\_11.variants.two\_pow\_two} in formalised-conjectures
repositories, together with a remark, which we have not independently
verified against a primary source, that Erd\H{o}s thought a two-power
relaxation of \#11 ``might be easy''. We give a sieve reduction in which the
dominant term is controlled \emph{exactly and uniformly in $n$} by a new
cancellation lemma at non-Wieferich primes (Lemma~\ref{lem:K}, the $p$-adic
reduction underlying the square-sieve literature, here specialised to this sieve),
close the boundary term completely (\S\ref{sec:R1}), reduce the remaining
worst-case obstruction to a single, clean, $n$-independent multiplicity
statement (Lemma~\ref{lem:M}, \S\ref{sec:M}), and prove that the two-powers
variant holds \emph{unconditionally for almost all $n$}
(Corollary~\ref{cor:almostall}, density $1-O(1/(L\log L))$, via a second moment on
$(L,L^2]$ and a first moment at a raised threshold on the wide range
$(L^2,\sqrt{n/2}]$).
For the \emph{every}-$n$ statement we give a conditional theorem
(Theorem~\ref{thm:conditional}): it holds provided a single grand-sieve bound
on the multiplicity distribution of $\{2^l+2^m\bmod p^2\}$ (hypothesis (TB)),
together with two softer structural inputs. Along the way we prove unconditionally
that the relevant order sum $\sum_{p\in(L,L^2]}\lfloor L/\ord_p(2)\rfloor$ is
$o(L^2)$ (Proposition~\ref{prop:ordersum}) and that the Type~A count satisfies
$\sum M_p=O(L^{8/3}/\log L)$ via García--Voloch (Proposition~\ref{prop:gv}),
reducing the obstruction to the Type~B tail. For that tail we show that it is
closed by \emph{circularity}: Bourgain--Chang gives
the additive sup-norm unconditionally (Theorem~\ref{thm:bc}), but the moment bound
needed to convert it into (TB) is \emph{equal} to the additive energy, hence to
(TB) itself (Proposition~\ref{prop:circ}). A function-field model (\S\ref{sec:ff})
locates the difficulty as \emph{archimedean} -- the failure of $\{2^j\}$ to be
exactly Sidon, powers carrying where monomials do not -- rather than a lack of
cancellation, and identifies (M$''$) at $k=2$ as a squarefree-value statement, a
cousin of \#11. The remaining gaps are stated honestly with the exact point where
elementary counting fails (\S\ref{sec:open}, \S\ref{sec:ff}). All are decoupled
from two classical obstructions one might expect to inherit: the
infinitude of non-Wieferich primes (only the two \emph{known} Wieferich
primes are ever excluded, an $O(1)$ cost), and Crocker's 1971 disproof of the
\emph{prime}-kernel analogue $n=p+2^a+2^b$, which does not transfer to a
squarefree kernel.
\end{abstract}

\tableofcontents

% ============================================================================
\section{Introduction}
% ============================================================================

\subsection{The problem and its variant}

Erd\H{o}s Problem~\#11 \cite{erdosproblems11} asks whether every sufficiently
large $n$ can be written as $n=k+2^m$ with $k$ squarefree and $m\ge0$. The
problem is open. Partial and conditional results exist in the
literature attributed (see the exact reference list of Kalinin
\cite{Kalinin2026}, who states and uses them) to Granville--Soundararajan,
Hercher, Crocker, and Platt--Trudgian; we rely on \cite{Kalinin2026} for the
precise statements and do not reproduce them here, since cross-checking each
primary source individually is outside the scope of this note.

We study the following relaxation, which lets $n$ be written as a squarefree
kernel plus \emph{two} powers of $2$ instead of one:

\begin{quote}\itshape
\textbf{Two-powers variant.} For every odd integer $n>1$ there exist
$k,l,m\ge0$ with $k$ squarefree and $n=k+2^l+2^m$.
\end{quote}

(Restricting to odd $n$ costs nothing: if $n$ is even, $n=k+2^l+2^m$ with $k$
odd forces $l=m=0$ to be excluded or absorbed trivially by parity, and the
even case reduces to a one-power-type statement on $n/2$-scale considerations
that we do not pursue here; all of our results concern odd $n$, matching the
Lean statement.) This statement is recorded as
\texttt{erdos\_11.variants.two\_pow\_two} in formalised-conjectures Lean
repositories. It is, to our knowledge, unpublished and unproved; this note
does not close it, but reduces it as far as we have been able to push it with
elementary and semi-elementary tools.

\subsection{Contributions}

\begin{enumerate}
\item \textbf{Lemma K} (\S\ref{sec:K}): an exact character-sum cancellation,
$\sum_{x\in\langle2\rangle}e_{p^2}(ax)=0$ for every non-Wieferich prime $p\ge3$
and every $a$ with $p\nmid a$, where $\langle2\rangle$ is the subgroup of
$(\Z/p^2)^\times$ generated by $2$. The vanishing mechanism is standard (the
$p$-adic reduction underlying the square-sieve literature); the contribution is
its exact specialisation to this sieve, which kills, \emph{exactly and uniformly
in $n$}, the dominant ``complete period'' term of every medium-sized prime with
no equidistribution input.
\item A complete elementary closure of the boundary term left over by
Lemma~K (Proposition~\ref{prop:R1}, \S\ref{sec:R1}): $O(L^2/\log L)=o(L^2)$,
where $L=\lfloor\log_2n\rfloor$.
\item A deterministic, $n$-independent reduction (Lemma~\ref{lem:M},
\S\ref{sec:M}) \emph{dissolving the worst-case-$n$ obstruction} on the single
dominant remaining range, to a clean additive multiplicity question.
\item \textbf{Type~A, unconditionally} (Proposition~\ref{prop:gv}): a bound of
García--Voloch on additive representations in the subgroup $\langle2\rangle\bmod
p$ gives $\sum_{\text{Type A}}M_p=O(L^{8/3}/\log L)$ unconditionally -- short of
the target $o(L^2)$ by exactly $L^{2/3}$.
\item \textbf{Type~B is closed by circularity} (\S\ref{sec:bc}): the
Bourgain--Chang additive sup-norm bound holds unconditionally on our sums
(Theorem~\ref{thm:bc}), but the moment bound that would convert it into the
Type~B hypothesis (TB) equals the additive energy, i.e.\ (TB) \emph{itself}
(Proposition~\ref{prop:circ}). Thus (TB) is a genuine open estimate, not a
statement awaiting the right reference.
\item \textbf{A function-field model} (\S\ref{sec:ff}) proves the $\F_q[t]$
analogue of (TB) by a degree count (Proposition~\ref{prop:ff}) and shows the
mechanism does not transpose, diagnosing the obstruction over $\Z$ as
\emph{archimedean} -- the failure of $\{2^j\}$ to be exactly Sidon.
\item An unconditional theorem (\S\ref{sec:almost}), via a second moment and
Markov, for \emph{almost all $n$} (density $1-O(1/(L\log L))$).
\item \textbf{Machine verification}: the elementary lemmas underpinning the
unconditional results are formally verified in Lean~4 against Mathlib
(\texttt{Erdos11\_verified.lean}, five lemmas, \texttt{0 sorry}).
\item A precise, honest statement of the two remaining gaps
(\S\ref{sec:open}, \S\ref{sec:ff}), each pinned to the exact point where
brute-force counting fails and at which scale.
\end{enumerate}

\subsection{Decoupling from two classical obstructions}\label{sec:intro-decoupling}

\paragraph{Independence from Wieferich-prime infinitude.} Lemma~K requires
$p$ to be non-Wieferich. It is open whether infinitely many Wieferich primes
exist (equivalently, whether infinitely many primes fail to be
non-Wieferich); only two are known, $1093$ and $3511$. Our reduction does
\emph{not} require the (open, and widely believed) statement that
non-Wieferich primes have density $1$, or even that there are infinitely
many: it only ever needs to discard the \emph{known} Wieferich primes
$\le\sqrt n$, an $O(1)$-many exclusion contributing $O(L)$ to every estimate
below, completely harmless. If new Wieferich primes are discovered, the
argument is unaffected; if it turned out that a positive density of primes
were Wieferich (which would be a sensational and unexpected
result), Lemma~K would have to be reproved for those primes by other means,
but this is not on the table.

\paragraph{Independence from Crocker's disproof.} Crocker \cite{Crocker1971}
disproved the \emph{prime}-kernel two-powers statement: there exist
infinitely many $n$ with $n\ne p+2^a+2^b$ for every prime $p$. His method is a
covering-congruence argument using Fermat numbers $2^{2^s}+1$: he exhibits
residue classes modulo a suitable Fermat number in which $2^a+2^b$ is forced
into a fixed congruence class that makes $n-2^a-2^b$ composite for every
choice of $a,b$ in the relevant range. This disproof does \emph{not} transfer
to the squarefree-kernel variant studied here, for a structural reason:
``composite'' is forced by a congruence condition of density essentially $1$
in the relevant residue classes (almost every integer in a Fermat-number
residue class is composite), whereas ``non-squarefree'' is a much rarer
condition (density $1-6/\pi^2\approx0.392$ overall, and in particular a
positive proportion of any fixed congruence class is squarefree, by
elementary sieve estimates). Crocker's congruence trick gives no control over
squarefreeness at all. This is one reason the squarefree-kernel variant is
plausible while its prime-kernel cousin is false.

\subsection{A structural cousin: Erd\H{o}s Problem~\#9}

Erd\H{o}s Problem~\#9 concerns $n=p+2^k+2^l$ with $p$ \emph{prime} (not
squarefree) -- the same surface shape (a fixed-size kernel condition plus two
powers of $2$) but a different, and harder, kernel condition. The toolkit
that controls a squarefree kernel (counting square divisors via a sieve, as
here) is structurally different from the toolkit needed for a prime kernel
(which needs control of prime gaps / sieve lower bounds, a strictly harder
problem). We do not address \#9 here; we mention it only to locate
\#11's two-powers variant precisely among nearby open problems of the same
shape.

% ============================================================================
\section{Setup and notation}\label{sec:setup}
% ============================================================================

Throughout, $n$ is odd, $n>1$, and $L:=\lfloor\log_2n\rfloor$. We count
\emph{ordered} pairs $(l,m)\in\{0,\dots,L\}^2$ (using ordered pairs only
inflates every count by a harmless factor close to $2$); there are
$T:=(L+1)^2$ of them. We discard the (negligibly many, $O(L)$) pairs with
$2^l+2^m\ge n$, since they cannot contribute a valid decomposition; this only
shrinks the pool of candidate pairs and makes every bound below
conservative.

\begin{definition}
For a prime $p\ge3$,
\[
N_p(n):=\#\{(l,m)\in\{0,\dots,L\}^2:\ p^2\mid n-2^l-2^m\}
=\#\{(l,m):\ 2^l+2^m\equiv n\pmod{p^2}\}.
\]
Write $\dpp:=\ord_{p^2}(2)$, $\ep:=\ord_p(2)$,
$\langle2\rangle:=\{2^l\bmod p^2\}$ (the cyclic subgroup of
$(\Z/p^2)^\times$ generated by $2$, of order $\dpp$), and
$H_p:=\{2^l\bmod p\}$ (the cyclic subgroup of $(\Z/p)^\times$ generated by
$2$, of order $\ep$). A prime $p$ is \emph{Wieferich} if $2^{p-1}\equiv1
\pmod{p^2}$; equivalently $\dpp=\ep$ (no extra factor of $p$ is gained from
$\Z/p$ to $\Z/p^2$). Only two Wieferich primes are known: $1093$ and $3511$.
For non-Wieferich $p$, $\dpp=p\cdot\ep$.
\end{definition}

Since $k=n-2^l-2^m$ is always odd, $p=2$ never occurs in any of the sieve
sums below.

\begin{proposition}[Basic reduction]\label{prop:star}
The number of pairs $(l,m)\in\{0,\dots,L\}^2$ for which $k=n-2^l-2^m$ is
\emph{not} squarefree is at most $\sum_{3\le p\le\sqrt n}N_p(n)$. If this sum
is $<T$, then the two-powers variant holds for $n$.
\end{proposition}
\begin{proof}
Union bound over primes $p$ with $p^2\mid k$ for some valid pair.
\end{proof}

(Small $n$ are not a concern: the variant has been checked directly, without
exception, up to $5\times10^7$.)

\begin{proposition}[Periodicity bound]\label{prop:periodic}
For every prime $p\ge3$, $\displaystyle N_p(n)\le\frac{(L+1)^2}{\dpp}+(L+1)$.
\end{proposition}
\begin{proof}
$2^l\bmod p^2$ is $\dpp$-periodic in $l$. For each fixed $l\le L$, the
equation $2^m\equiv n-2^l\pmod{p^2}$ has at most
$\lceil(L+1)/\dpp\rceil\le(L+1)/\dpp+1$ solutions $m\in\{0,\dots,L\}$. Summing
over the $L+1$ values of $l$ gives the claim.
\end{proof}

\subsection{The pivot constant}\label{sec:pivot}

The reduction below needs $\sum_{p\ge3}1/\dpp<\tfrac12$. We record two facts:
the precise numerical value, and an elementary proof that the underlying
infinite series in fact converges (so that summing it over an arbitrarily
large range of primes, as we will, is legitimate).

\begin{proposition}[Convergence]\label{prop:convergence}
For every prime $p\ge3$, $\ep>\log_2p-1$. Consequently, for every
non-Wieferich $p$, $\dpp=p\,\ep>p(\log_2p-1)$, and
$\sum_{p\ge3}1/\dpp$ converges (the two Wieferich primes contribute two
extra, harmless, finite terms).
\end{proposition}
\begin{proof}
By definition $p\mid2^{\ep}-1<2^{\ep}$, so $p<2^{\ep}$, i.e.\
$\ep>\log_2p$; we used the slightly weaker $\ep>\log_2p-1$ to avoid
worrying about the base case. For the convergence: by Mertens' classical
theorem, $\sum_{p\le x}1/p=\log\log x+M+o(1)$; partial summation against the
decreasing weight $1/(\log_2p-1)$ turns this into the classical fact that
$\sum_p\frac1{p\log p}$ converges (the density of primes near $t$ is
$1/\log t$ by the prime number theorem, so the partial sums behave like
$\int^x\frac{dt}{t(\log t)^2}$, convergent). Since $1/\dpp<1/(p(\log_2p-1))=
O(1/(p\log p))$ for non-Wieferich $p$, comparison gives convergence of
$\sum_p1/\dpp$.
\end{proof}

\begin{proposition}[Numerical value]\label{prop:pivot}
$\sum_{3\le p<5000}1/\dpp=0.320516\ldots$, computed exactly via exact
modular-order arithmetic (script \texttt{orders.py}), dominated by $p=3$
($1/6$), $p=5$ ($1/20$), $p=7$ ($1/21$). By Proposition~\ref{prop:convergence}
the tail $p\ge5000$ is provably summable and numerically negligible (a crude
comparison with $\sum_{p\ge5000}1/(p(\log_2p-1))$ already gives less than
$3\times10^{-4}$); in particular
\[
S:=\sum_{p\ge3}\frac1{\dpp}\ \approx\ 0.3205\ <\ \frac12.
\]
\end{proposition}

\begin{remark}
For the application below we only ever need a \emph{fixed} threshold $z$
(we use $z=5000$) with $\sum_{p<z}1/\dpp<\tfrac12$ exactly, plus the fact
that the tail $\sum_{p\ge z}1/\dpp\to0$ as $z\to\infty$ -- the latter is
exactly the convergence statement of Proposition~\ref{prop:convergence}, not
a numerical estimate.
\end{remark}

% ============================================================================
\section{Lemma K: exact cancellation at non-Wieferich primes}\label{sec:K}
% ============================================================================

The vanishing mechanism below -- a multiplicative subgroup whose order is
divisible by $p$ necessarily contains the kernel $K=1+p\Z/p^2\Z$ of the
reduction $(\Z/p^2)^\times\to(\Z/p)^\times$, and an additive character summed
over a union of $K$-cosets vanishes -- is standard in the theory of exponential
sums modulo prime powers (it is the elementary $p$-adic reduction underlying,
e.g., Cochrane--Zheng). What is specific to this note is only its exact
application to the present sieve; we state it as a lemma for that reason, not
because the identity itself is new.

\begin{lemma}[Lemma K]\label{lem:K}
Let $p\ge3$ be a non-Wieferich prime. For every integer $a$ with $p\nmid a$,
\[
S(a):=\sum_{x\in\langle2\rangle}e_{p^2}(ax)\ =\ 0,
\qquad e_{p^2}(t):=e^{2\pi it/p^2}.
\]
\end{lemma}
\begin{proof}
Since $p$ is non-Wieferich, $\dpp=p\,\ep$, so $\langle2\rangle$ (cyclic of
order $\dpp$) contains the unique subgroup of order $p$ of
$(\Z/p^2)^\times$, namely the kernel
$K=\{1+jp:\ j=0,\dots,p-1\}$ of the reduction map
$(\Z/p^2)^\times\to(\Z/p)^\times$. (Indeed $\langle2\rangle\to H_p$ is
surjective with kernel $\langle2\rangle\cap K$; comparing orders $p\,\ep$ and
$\ep$, this kernel has order $p$, hence equals $K$.) Partition
$\langle2\rangle$ into cosets $gK$ over a set of representatives $g$:
\[
S(a)=\sum_{g}\sum_{j=0}^{p-1}e_{p^2}\bigl(ag(1+jp)\bigr)
=\sum_{g}e_{p^2}(ag)\underbrace{\sum_{j=0}^{p-1}e_p(agj)}_{=0},
\]
the inner sum being geometric with ratio $e_p(ag)\ne1$ since $p\nmid a$ and
$p\nmid g$ (as $g\in(\Z/p^2)^\times$).
\end{proof}

\begin{remark}[Numerical verification]
For $p=3,5,7,11,13,17,19,23,29,31,37,41,43,127$ (script \texttt{lemmaE2.py}),
all non-Wieferich, the kernel $K$ is exactly contained in $\langle2\rangle$
and the maximum of $|S(a)|$ over $p\nmid a$ is exactly $0$, confirming the
lemma. For the two known Wieferich primes $1093$ and $3511$, $\dpp=\ep$
(no extra factor of $p$), $K\not\subseteq\langle2\rangle$, and $S(a)\ne0$:
the mechanism breaks exactly where the hypothesis fails, as expected.
\end{remark}

\subsection{Exact structure on a complete period}

\begin{proposition}[Exact period count]\label{prop:exactperiod}
Let $p$ be non-Wieferich. For every residue $n\bmod p^2$,
\[
\#\bigl\{(l,m)\in\{0,\dots,\dpp-1\}^2:\ 2^l+2^m\equiv n\pmod{p^2}\bigr\}
\ =\ p\cdot r_p(n),
\qquad
r_p(n):=\#\bigl\{(u,v)\in H_p^2:\ u+v\equiv n\pmod p\bigr\}.
\]
\end{proposition}
\begin{proof}[Proof sketch]
By Lemma~\ref{lem:K} (via Fourier inversion on $\langle2\rangle$, or
directly): each solution $(u,v)\in H_p^2$ of $u+v\equiv n\pmod p$ lifts to
exactly $p$ pairs $(x,y)\in\langle2\rangle^2$ with $x+y\equiv n\pmod{p^2}$,
namely the $p$ choices of $x$ in the $K$-coset of $u$ (each forcing $y$'s
residue mod $p$ to be $v$, hence $y$'s $K$-coset, but \emph{not} pinning $y$
further -- the count is exactly balanced by the geometric cancellation of
Lemma~\ref{lem:K}). Verified by direct computation for
$p=3,5,7,11,13,17,19,23$ (script \texttt{verify\_structure.py}): the identity
holds exactly for every residue $n\bmod p^2$.
\end{proof}

Since $r_p(n)\le|H_p|=\ep$ (each $u\in H_p$ admits at most one matching $v$),
this gives, for non-Wieferich $p$ with $\dpp\le L+1$ (so that
$Q_p:=\lfloor(L+1)/\dpp\rfloor\ge1$ complete periods fit in $\{0,\dots,L\}$):
\[
Q_p^2\,p\,r_p(n)\ \le\ Q_p^2\,p\,\ep\ =\ Q_p^2\,\dpp\ \le\ \frac{(L+1)^2}{\dpp},
\qquad\text{uniformly in } n.
\]

\begin{theorem}[Complete periods are uniformly $o(L^2)$]\label{thm:complete}
For every odd $n>1$,
\[
\sum_{\substack{p\ \mathrm{non\text{-}Wief.}\\ \dpp\le L+1}}
Q_p^2\,p\,r_p(n)\ \le\ (L+1)^2\sum_{p\ge3}\frac1{\dpp}\ \le\ 0.3205\,(L+1)^2,
\]
\emph{uniformly in $n$, with no equidistribution hypothesis whatsoever}.
\end{theorem}
\begin{proof}
Immediate from the displayed bound above, summed over the (at most
$\pi(L+1)$, see Proposition~\ref{prop:locate} below) primes with
$\dpp\le L+1$, and Proposition~\ref{prop:pivot}.
\end{proof}

This is the main term of the whole sieve, and it is now controlled
\emph{exactly}, for every $n$, by an elementary cancellation -- no
equidistribution, no analytic input, no conjecture.

% ============================================================================
\section{Closing the boundary term: Proposition R1}\label{sec:R1}
% ============================================================================

Theorem~\ref{thm:complete} controls the contribution of full periods. What
remains, for primes with $\dpp\le L+1$, is the boundary correction between a
whole number of periods and the actual finite range $\{0,\dots,L\}$.

\begin{proposition}[Localisation]\label{prop:locate}
If $p$ is non-Wieferich and $\dpp\le L+1$, then $p\le L+1$.
\end{proposition}
\begin{proof}
$\dpp=p\,\ep\ge p$ since $\ep\ge1$.
\end{proof}

So the primes carrying a boundary term are confined to $p\le L+1$, at most
$\pi(L+1)$ of them.

\begin{proposition}[R1, elementary closure]\label{prop:R1}
For $n$ odd, define
\[
R1(n):=\sum_{\substack{p\ \mathrm{non\text{-}Wief.}\\ \dpp\le L+1}}
\bigl(N_p(n)-Q_p^2\,p\,r_p(n)\bigr).
\]
Then $\displaystyle R1(n)\ <\ 3(L+1)\,\pi(L+1)\ =\ O\!\left(\frac{L^2}{\log L}\right)=o(L^2)$.
\end{proposition}
\begin{proof}
Fix such a $p$, and write $L+1=Q_pd_p+s_p$, $0\le s_p<d_p$. On the range
$\{0,\dots,L\}$, each residue $x\in\langle2\rangle$ is hit
$c(x)=Q_p+\varepsilon(x)$ times, $\varepsilon(x)\in\{0,1\}$, with
$\#\{x:\varepsilon(x)=1\}=s_p$ (the residues $2^0,\dots,2^{s_p-1}$). Writing
$P_n:=\{x\in\langle2\rangle:n-x\in\langle2\rangle\}$ (so $|P_n|=p\,r_p(n)$ by
Proposition~\ref{prop:exactperiod}),
\[
N_p(n)=\sum_{x\in P_n}c(x)c(n-x)
=Q_p^2|P_n|+2Q_p\gamma_p(n)+\delta_p(n),
\]
with $\gamma_p(n):=\#\{x\in P_n:\varepsilon(x)=1\}\le s_p$ and
$0\le\delta_p(n)\le\gamma_p(n)$. Hence
$N_p(n)-Q_p^2p\,r_p(n)=2Q_p\gamma_p(n)+\delta_p(n)\le(2Q_p+1)s_p$. Since
$Q_ps_p<Q_pd_p\le L+1$ and $s_p<d_p\le L+1$,
\[
(2Q_p+1)s_p=2Q_ps_p+s_p<2(L+1)+(L+1)=3(L+1).
\]
Summing over the (at most $\pi(L+1)$) primes with $\dpp\le L+1$ gives the
bound; the prime number theorem gives $\pi(L+1)\sim(L+1)/\log(L+1)$. The
$\le2$ known Wieferich primes contribute, separately and trivially, $O(L)$
each.
\end{proof}

\begin{remark}[Numerical confirmation]
Script \texttt{verify\_residual.py} computes $R1(n)$ exactly for the worst
$n$ at every dyadic scale up to $n\approx2^{22}$ ($L=13,\dots,21$). The
maximal prime ever carrying a boundary term in this range is $7$ (far below
the bound $L+1$), and the measured $R1\in\{0,4,8,\dots,32\}$ stays an order of
magnitude below the proved bound $3(L+1)\pi(L+1)\in[234,528]$ over the same
range.
\end{remark}

Combining Theorem~\ref{thm:complete} and Proposition~\ref{prop:R1}: for every
odd $n$, the total contribution of primes with $\dpp\le L+1$ (\emph{all}
primes other than the $\le2$ known Wieferich exceptions) is at most
\[
0.3205\,(L+1)^2+O\!\left(\frac{L^2}{\log L}\right)+O(L)
\ =\ \bigl(0.3205+o(1)\bigr)(L+1)^2,
\]
uniformly in $n$ -- an elementary, rigorous, complete closure leaving a
margin of $\approx0.18\,T$ before the threshold $T=(L+1)^2$ of
Proposition~\ref{prop:star}.

% ============================================================================
\section{Lemma M: deterministic dissolution of the worst case}\label{sec:M}
% ============================================================================

What remains is the contribution of primes with $\dpp>L+1$ (no complete
period inside $\{0,\dots,L\}$); by Proposition~\ref{prop:locate} (contrapositive)
this is exactly the set of primes $p>L+1$ that are non-Wieferich, together
with the (negligible) Wieferich exceptions. Write
\[
R2(n):=\sum_{\substack{p\le\sqrt n\\ \dpp>L+1}}N_p(n).
\]
By the trivial bound $N_p(n)\le L+1$ (distinct $2^l\bmod p^2$, no period),
combined with the prime-counting bound $\#\{p\le\sqrt n\}\sim\sqrt n/\log\sqrt n$,
elementary counting alone gives only $R2(n)\le(L+1)\sqrt n$, hopelessly larger
than $T=(L+1)^2$. The empirical data (\cite{erdos11data}) shows $R2(n)$ is in
fact always $\ll L^2$ and concentrated on a bounded range of primes just
above $L$; this section isolates, deterministically and without reference to
any particular $n$, exactly how much of this is provable.

We restrict attention to the range $p\in(L,L^2]$, which the data identifies
as carrying the dominant share of $R2$; the remaining two sub-ranges of $R2$
are discussed in \S\ref{sec:open}.

\begin{definition}
For a prime $p$, set $M_p:=\max_{n}N_p(n)$ ($n$ ranging over odd integers,
with $L$ implicitly determined by $n$ -- equivalently, $M_p$ is the maximal
number of representations of any residue $n\bmod p^2$ as $2^l+2^m$ with
$0\le l,m\le L$, for the relevant $L$). This is an $n$-independent quantity.
\end{definition}

\begin{lemma}[Lemma M -- deterministic dissolution]\label{lem:M}
For \emph{every} odd $n$,
\[
\sum_{p\in(L,L^2]}N_p(n)\ \le\ \sum_{p\in(L,L^2]}M_p.
\]
\end{lemma}
\begin{proof}
$N_p(n)\le\max_{n'}N_p(n')=M_p$ for every $p$, by definition; sum over $p$.
\end{proof}

This trivial-looking statement is the key structural gain of this section:
\emph{it eliminates the worst-case-$n$ problem entirely} for this range. If
$\sum_{p\in(L,L^2]}M_p=o(L^2)$, this range contributes $o(L^2)$ to
\emph{every} $n$ simultaneously, deterministically -- no averaging, no
exceptional set.

\subsection{The unconditional half: the Sidon decomposition}

If $2^i+2^j\equiv2^{i'}+2^{j'}\pmod{p^2}$ has no solution with
$\{i,j\}\ne\{i',j'\}$ (both pairs ranging over $\{0,\dots,L\}$, $L$ ranging
implicitly over all values relevant to $M_p$), we call $p$ \emph{Sidon}; then
every value $n\bmod p^2$ has at most one unordered representation, so the
ordered multiplicity is at most $2$ (the pair and its swap).

\begin{proposition}[Unconditional Sidon term]\label{prop:sidonterm}
\[
\sum_{p\in(L,L^2]}M_p\ =\ \underbrace{2\bigl(\pi(L^2)-\pi(L)\bigr)}_{\text{Sidon contribution}}
\ +\ \underbrace{\sum_{p\in(L,L^2]\,\mathrm{non\text{-}Sidon}}(M_p-2)}_{\text{excess}},
\]
and the Sidon contribution alone satisfies
$2(\pi(L^2)-\pi(L))=O(L^2/\log L)=o(L^2)$, \emph{unconditionally} (immediate
from the prime number theorem; no structural input about $\{2^l\bmod p^2\}$
is needed for this half).
\end{proposition}
\begin{proof}
For Sidon $p$, $M_p=2$ exactly (a representation and its swap; $M_p\ge2$
trivially once $L\ge1$). The displayed identity is then a rearrangement of
the sum. The bound on the Sidon term is the prime-counting estimate
$\pi(L^2)-\pi(L)\sim L^2/(2\log L)$.
\end{proof}

The remaining \emph{excess} term -- the total multiplicity overshoot summed
over non-Sidon primes in $(L,L^2]$ -- is exactly the open quantity discussed
in \S\ref{sec:open} (named M$''$ there). Brute-force counting of it already
fails at the first step: the crude bound $M_p-2\le L-1$ (trivial, since
$M_p\le L+1$) combined with the trivial bound on the number of non-Sidon
primes ($\le\pi(L^2)-\pi(L)$) gives only $\Theta(L^2)$, not $o(L^2)$ -- the
excess is genuinely smaller than this in the data (by roughly a factor
$1/\log L$), but proving it requires controlling the \emph{distribution} of
the multiplicities $M_p$, not just their existence.

\subsection{The order sum is unconditionally $o(L^2)$}\label{sec:ordersum}

We isolate one half of the excess that \emph{is} unconditional. Split the range
by the multiplicative order of $2$: call $p\in(L,L^2]$ \emph{Type~A} if $\ep\le L$
and \emph{Type~B} if $\ep>L$.

\begin{proposition}[Order sum]\label{prop:ordersum}
$\displaystyle\sum_{p\in(L,L^2]}\Bigl\lfloor\tfrac{L}{\ep}\Bigr\rfloor
=O\!\left(\tfrac{L^2}{\log L}\right)=o(L^2)$.
\end{proposition}
\begin{proof}
$\lfloor L/d\rfloor=\#\{j\in[1,L]:d\mid j\}$, and $\ep\mid j\iff p\mid2^j-1$, so
$\sum_{p}\lfloor L/\ep\rfloor=\sum_{j\le L}\#\{p\in(L,L^2]:p\mid2^j-1\}$. The
primes $p>L$ dividing $2^j-1$ have product $\le2^j-1<2^j$; if there are $m$ of
them then $L^m<2^j$, so $m<j\log2/\log L$. Summing,
$\sum_{j\le L}j\log2/\log L=(\log2/\log L)\cdot L(L+1)/2=O(L^2/\log L)$.
\end{proof}

The gain over the trivial bound (which counts each $2^j-1$ with its full $O(j)$
prime factors and gives $\Theta(L^2)$) is a factor $\log L$, and it comes
\emph{entirely} from the constraint $p>L$: large primes cannot divide $2^j-1$
too often. This is the only place the lower endpoint of the range is used.

\subsection{Classification of $(L,L^2]$: the involution structure}\label{sec:classify}

Each pair contributes $2^l+2^m=2^l(1+2^{m-l})$; a fibre's arithmetic is governed
by the \emph{gap values} $1+2^\delta$ ($\delta=|m-l|$), which satisfy an exact
identity.

\begin{lemma}[Involution]\label{lem:involution}
For every $\delta$, $\ 1+2^{\,\ep-\delta}\equiv2^{-\delta}(1+2^\delta)\pmod p$.
\end{lemma}
\begin{proof}
$2^{-\delta}(1+2^\delta)=2^{-\delta}+1=2^{\,\ep-\delta}+1$, using $2^{\ep}\equiv1
\pmod p$.
\end{proof}

Thus $\delta$ and $\ep-\delta$ send $1+2^\delta$ into the same coset of
$\langle2\rangle\bmod p$: the gap values pair up under $\delta\mapsto\ep-\delta$.
Using the $p$-adic refinement $2^{\ep}\equiv1+w_pp\pmod{p^2}$ (with $w_p\ne0$
exactly for non-Wieferich $p$), the representations of a fixed residue $n\bmod
p^2$ organise into arithmetic progressions in $\delta$ of common difference a
multiple of $\ep$, giving a per-class bound.

\begin{proposition}[Per-class bound]\label{prop:perclass}
Let $p\in(L,L^2]$ be non-Wieferich. For any residue $n\bmod p^2$, the pairs
$(l,m)\in\{0,\dots,L\}^2$ with $2^l+2^m\equiv n$, grouped by $\delta\bmod\ep$,
form at most $\lfloor L/\ep\rfloor+1$ pairs per class. Hence
\[
M_p\ \le\ \kappa_p\,(\lfloor L/\ep\rfloor+1),\qquad
\kappa_p:=\#\{\text{active classes }\delta\bmod\ep\text{ in a maximal fibre}\}.
\]
\end{proposition}
\begin{proof}[Proof sketch]
Since $\dpp=p\,\ep>L$, for each $l$ there is at most one $m\in\{0,\dots,L\}$ with
$2^m\equiv n-2^l\pmod{p^2}$. Within a class $(l\bmod\ep,\delta\bmod\ep)$, writing
$l=l_0+\ep a$, $\delta=\delta_0+\ep b$, the congruence reduces (via
$2^{\ep}\equiv1+w_pp$, $w_p\ne0$) to a single linear congruence
$aX+bY\equiv Z\pmod p$ in the triangle $\{a,b\ge0,\ a+b\le(L-l_0-\delta_0)/\ep\}$
of side $T\le L/\ep<p$; as $T<p$, each $b$ admits at most one $a$, giving
$\le\lfloor L/\ep\rfloor+1$ pairs. Sum over the $\kappa_p$ active classes.
(Verified with no exceptions for all $p\in(L,L^2]$, $L\le160$; non-degeneracy of
the linear form is exactly the non-Wieferich condition.)
\end{proof}

Empirically $\kappa_p\le2$ for \emph{genuinely small} order ($\ep\ll L$: the
active classes form at most one involution orbit plus a fixed point), whereas for
\emph{boundary} primes $\ep\approx L$ (where $2$ is near-primitive) $\kappa_p$
behaves like a Poisson maximum and grows slowly (measured $\le6$ for $L\le200$).
We do not prove a bound; we isolate as hypothesis (TA) the mild growth condition
$\max_{\ep\le L}\kappa_p=o(\log L)$. Since $\sum_p\lfloor L/\ep\rfloor$ already
carries a $1/\log L$ saving (Proposition~\ref{prop:ordersum}), this suffices:
\[
\sum_{p\in(L,L^2]:\,\ep\le L}M_p\ \le\ \Bigl(\max_{\ep\le L}\kappa_p\Bigr)
\Bigl(\sum_p\lfloor L/\ep\rfloor+\pi(L^2)\Bigr)
\ =\ o(\log L)\cdot O\!\Bigl(\tfrac{L^2}{\log L}\Bigr)\ =\ o(L^2)\quad\text{[under (TA)].}
\]

\emph{A partial unconditional bound on $\kappa_p$.} The empirical $\kappa_p\le2$
is not needed in full: a theorem of García--Voloch on additive representations in
a multiplicative subgroup gives a non-trivial bound directly.

\begin{proposition}[Unconditional $\kappa_p\ll\ep^{2/3}$, via García--Voloch]\label{prop:gv}
For $G=\langle2\rangle\bmod p$ with $|G|=\ep<(p-1)/((p-1)^{1/4}+1)$ (so $\ep\lesssim
p^{3/4}$) and any $b\not\equiv0\pmod p$, one has $r_p(b):=\#\{(u,v)\in G^2:u+v\equiv b\}
\le4\,\ep^{2/3}$ \textup{(}García--Voloch~\cite{GarciaVoloch1988}, in the exact
form of Konyagin's Theorem~2.1~\cite{Konyagin-notes}; Heath-Brown gave a
Stepanov-method proof, and Heath-Brown--Konyagin the energy $T_2(G)\ll\ep^{5/2}$ for
$\ep\le p^{2/3}$\textup{)}. Reducing a maximal fibre modulo $p$ sends its $\kappa_p$
active classes to distinct solutions of $u+v\equiv n\bmod p$ with $u,v\in G$ (two
pairs in different classes $\delta\bmod\ep$ give different gap values $1+2^\delta$,
hence different representations mod $p$); so $\kappa_p\le r_p(n\bmod p)\ll\ep^{2/3}$
unconditionally, whence
\[
\sum_{p\in(L,L^2]:\,\ep\le L}M_p\ =\ O\!\left(\frac{L^{8/3}}{\log L}\right).
\]
\end{proposition}
\begin{proof}
With $M_p\le\kappa_p(\lfloor L/\ep\rfloor+1)$ and $\#\{p>L:\ep=e\}\ll e/\log L$
(the counting behind Proposition~\ref{prop:convergence}),
$\sum M_p\ll(\log L)^{-1}\sum_{e\le L}e^{2/3}(L+e)\asymp L^{8/3}/\log L$. Two ranges
use the trivial $M_p\le L+1$ instead. The size hypothesis fails only for $\ep>p^{3/4}$,
i.e.\ $p<\ep^{4/3}\le L^{4/3}$: a band of $\le\pi(L^{4/3})=O(L^{4/3}/\log L)$ primes,
total $O(L^{7/3}/\log L)$. The excluded residues $b\equiv0\pmod p$ occur only when
$p\mid n$, for $\le\log n/\log L=O(L/\log L)$ primes, total $O(L^2/\log L)$; the
$p^2$-degenerate residue ($2^l+2^m\equiv0\bmod p^2$) is unreachable in range for
Type~A, as it needs $m-l\asymp\ord_{p^2}(2)/2>L$. Both bands are $o(L^{8/3}/\log L)$,
absorbed. \textup{(}Numerically the band holds $8,15,16,23,43$ primes for
$L=60,\dots,300$, always $\max p<L^{4/3}$, and the ratio $r_p(b)/\ep^{2/3}<1$ with no
drift to $\ep\approx900$; scripts \texttt{chantierD\_stepanov.py},
\texttt{chantierG\_band.py}.\textup{)}
\end{proof}

This is short of the target $o(L^2)$ by a factor $L^{2/3}$. Hypothesis (TA) is thus
now the sharper open question of whether the exponent $2/3$ can be lowered to $o(1)$
on the special fibres $1+2^\delta$ -- a rigidity that the generic Stepanov argument
ignores but the involution orbit structure may supply.

\subsection{The Type~B residual as a factorial-moment condition}\label{sec:typeB}

For Type~B primes ($\ep>L$) the fibres carry no progression
($\lfloor L/\ep\rfloor=0$) and $M_p$ is governed by coincidences alone. The sharp
reduction is through factorial moments. For $k\ge2$ put
\[
E_k^{\mathrm{tot}}:=\!\!\sum_{\substack{p\in(L,L^2]\\ \ep>L}}\sum_r\binom{N_p(r)}{k},
\qquad
E_k^{\mathrm{null}}:=\!\!\sum_{\substack{p\in(L,L^2]\\ \ep>L}}\frac{\binom{P}{k}}{(p^2)^{k-1}},
\quad P:=\tbinom{L+2}{2},
\]
$E_k^{\mathrm{null}}$ being the $k$-th factorial moment of the multinomial ($P$
balls in $p^2$ bins) summed over the range -- the ``random'' benchmark.

\begin{lemma}[Factorial-moment tail]\label{lem:markov}
$\#\{p\in(L,L^2]:\ep>L,\ M_p\ge k\}\le E_k^{\mathrm{tot}}$, and
$E_k^{\mathrm{null}}\sim L^3/\bigl(2^kk!\,(2k-3)\log L\bigr)$, which crosses $1$ at
$k^\ast\sim3\log L/\log\log L$.
\end{lemma}
\begin{proof}
Each $r$ with $N_p(r)\ge k$ contributes $\binom{N_p(r)}{k}\ge1$, and each such
$p$ has at least one such $r$; combine. The asymptotic uses
$\binom{P}{k}\sim P^k/k!$, $P\sim L^2/2$, and
$\sum_{p>L}p^{-(2k-2)}\sim L^{-(2k-3)}/((2k-3)\log L)$.
\end{proof}

\begin{theorem}[Conditional every-$n$ variant]\label{thm:conditional}
Suppose:
\begin{itemize}
\item[\textup{(TB)}] there is an absolute $C$ with
$E_k^{\mathrm{tot}}\le C^kE_k^{\mathrm{null}}$ for all $k\ge2$; only the tail
$k>k^\ast=\lfloor3\log L/\log\log L\rfloor$ is actually used, the range
$2\le k\le k^\ast$ being unconditional (see the proof);
\item[\textup{(TA)}] $\displaystyle\max_{p\in(L,L^2]:\,\ep\le L}\kappa_p=o(\log L)$;
\item[\textup{(WB)}] $\sum_{p\in(L^2,\sqrt{n/2}]}N_p(n)=o(L^2)$ uniformly in $n$.
\end{itemize}
Then the two-powers variant holds for every odd $n>n_0$.
\end{theorem}
\begin{proof}
By Proposition~\ref{prop:star} it suffices that $\sum_{3\le p\le\sqrt n}N_p(n)<T$.
The range $p\le L+1$ contributes $(0.3205+o(1))T$ (Theorem~\ref{thm:complete},
Proposition~\ref{prop:R1}); $(\sqrt{n/2},\sqrt n]$ contributes $O(L)$
(Kalinin~\cite{Kalinin2026}); $(L^2,\sqrt{n/2}]$ contributes $o(L^2)$ by (WB). For
$(L,L^2]$, by Lemma~\ref{lem:M} it suffices that $\sum_{p\in(L,L^2]}M_p=o(L^2)$;
Type~A is $o(L^2)$ under (TA) by \S\ref{sec:classify}, and Type~B is, under (TB)
and Lemma~\ref{lem:markov},
\[
\sum_{\ep>L}(M_p-1)=\sum_{k\ge2}\#\{M_p\ge k\}
\ \le\ \underbrace{k^\ast\,\pi(L^2)}_{\text{unconditional}}
+\underbrace{\sum_{k>k^\ast}\!C^kE_k^{\mathrm{null}}}_{\text{needs (TB), }k>k^\ast}
\ =\ O\!\Bigl(\tfrac{L^2}{\log\log L}\Bigr)=o(L^2).
\]
The first term is unconditional: each $\#\{M_p\ge k\}\le\pi(L^2)$, and there are
$\le k^\ast=O(\log L/\log\log L)$ of them, giving $O(L^2/\log\log L)$. Only the tail
$k>k^\ast$ invokes (TB), and there $\sum_{k>k^\ast}E_k^{\mathrm{null}}=o(1)$ since
$E_k^{\mathrm{null}}$ has already crossed $1$ at $k^\ast$. Summing the ranges gives
$(0.3205+o(1))T+o(L^2)<T$ for $n$ large.
\end{proof}

Hypothesis (TB) is the single genuinely analytic input, and it is equivalent to a
\emph{square-sieve / large-sieve} bound. Since $\{2^k\}$ is a Sidon set over
$\Z$, $E_k^{\mathrm{tot}}$ counts $k$-fold square-divisor incidences
$\#\{p>L:p^2\mid\sum_i\varepsilon_i2^{j_i}\}$ of power combinations; the natural
tool is Heath-Brown's square sieve~\cite{HeathBrown1984}. Its main term
reproduces $E_k^{\mathrm{null}}$ (the correct order, obtained by averaging over
$p$: per-prime bounds are a factor $L$ too weak), and its error term is a
\emph{mean-square} bound on the incomplete character sums $\sum_{j\le L}\chi(2^j)$
over $\{2^j\bmod p^2\}$. These are lacunary, so on the torus
$\|\sum a_je(2^j\theta)\|_{2k}\ll\sqrt k\,\|a\|_2$
(Rudin~\cite{Rudin1960}); the corresponding $\bmod\,p^2$ bound holds empirically
but reaches size $\sim L$ (not $O(1)$) on quasi-trivial characters, so a rigorous
proof needs the full large sieve, not a pointwise estimate. This is the same
incomplete-exponential-sum core as the fourth toolkit of \S\ref{sec:open}, now in
mean-square rather than sup-norm form. Empirically (TB) holds with margin:

\begin{center}
\begin{tabular}{lccc}
\toprule
$L$ & $E_2^{\mathrm{tot}}/E_2^{\mathrm{null}}$ & $E_3^{\mathrm{tot}}/E_3^{\mathrm{null}}$ & $E_4^{\mathrm{tot}}/E_4^{\mathrm{null}}$\\
\midrule
$80$ & $0.91$ & $0.91$ & $0.40$\\
$120$ & $0.95$ & $0.70$ & $0.35$\\
$160$ & $0.95$ & $0.86$ & $0.67$\\
\bottomrule
\end{tabular}
\end{center}

All ratios are $\le1$: the arithmetic fibres are \emph{sub}-multinomial at every
order, and $E_k^{\mathrm{tot}}$ already vanishes by $k=5$ at these $L$ (well below
the multinomial prediction). Restricting to Type~B removes the
super-concentrated Type~A fibres, whose energy is order-driven, not random. Note
that no independence or negative-association hypothesis is used: the passage from
the moments to the maximum is the elementary first-moment (Markov) bound of
Lemma~\ref{lem:markov}.

\begin{remark}[Which moment is binding]\label{rem:binding}
It is worth stating precisely which part of (TB) is used, to avoid a common
misreading. The Type~B contribution is
$\sum_{\ep>L}(M_p-1)=\sum_{k\ge2}\#\{p:M_p\ge k\}$, and \emph{every} term is
bounded trivially by $\#\{p:M_p\ge k\}\le\pi(L^2)=o(L^2)$; in particular the
$k=2$ term is free, and the individual energy $E_2^{\mathrm{tot}}=\sum_p\Delta_p$
is \emph{not} $o(L^2)$ (it is $\Theta(L^3/\log L)$, and need not be small). What
must be shown is only that the sum has $O(\pi(L^2))$ nonzero terms, i.e.\ that the
\emph{maximum} multiplicity satisfies $\max_{\ep>L}M_p=o(\log L)$; then
$M''_B\le(\max M_p)\,\pi(L^2)=o(L^2)$. Thus the binding instance of (TB) is the
\emph{large}-$k$ one, $k\asymp\log L/\log\log L$ (a large-deviation statement about
the maximum), not the variance $k=2$. The $k=2$ case is merely the cleanest window
onto the mechanism forcing sub-Poisson behaviour, namely that $\{2^k\}$ is Sidon
over $\Z$, so a collision mod $p^2$ forces $p^2\mid N$ with $N\ne0$, an event of
density $\sum_p1/p^2$ (verified: the density heuristic reproduces
$E_2^{\mathrm{tot}}$ to within $5\%$, stably in $L$). This mechanism is
base-independent: the same holds for $\{g^k\}$, any $g\ge2$.
\end{remark}

\subsection{The additive sup-norm, and why (TB) is closed by circularity}\label{sec:bc}

The square-sieve error term above is phrased through \emph{multiplicative}
characters $\chi$ of $(\Z/p^2)^\times$, and there $\sum_{j\le L}\chi(2^j)$ is a
geometric progression that reaches size $\sim L$ when $\chi$ is quasi-trivial on
$\langle2\rangle$ -- no cancellation. The \emph{additive} side behaves better,
and unconditionally so. Write $S(\xi):=\sum_{j\le L}e_{p^2}(\xi2^j)$.

\begin{theorem}[Additive sup-norm, Bourgain--Chang~\cite{BourgainChang2006}, Cor.~4.5]\label{thm:bc}
Let $q=\prod p_\alpha^{\nu_\alpha}$ have a bounded number of prime factors and
$\theta\in\Z_q^\times$ with $\ord_{p}(\theta)>q^\delta$ for every $p\mid q$. Then
for $t>q^\delta$, $\ \max_{\xi\in\Z_q^\times}|\sum_{s\le t}e_q(\xi\theta^s)|<t^{1-\varepsilon}$
with $\varepsilon=\varepsilon(\delta)>0$ absolute. For $q=p^2$, $\theta=2$,
$t=L+1$ and $p\in(L,L^2]$ Type~B, take $\delta=\tfrac15$: then $\ep>L\ge q^{1/5}$
and $L+1>q^{1/5}$, so \[|S(\xi)|<(L+1)^{1-\varepsilon}\quad(p\nmid\xi),\qquad
|S(p\eta)|<(L+1)^{1-\varepsilon'}\ (\eta\not\equiv0),\] the latter by reduction
$\bmod\,p$ \textup{(}order $\ep>\sqrt p$\textup{)}.
\end{theorem}

Numerically $\max_{p\nmid\xi}|S(\xi)|/(L+1)=0.37,0.31,0.24$ at $L=40,80,120$ on
Type~B primes -- well below $1$ and decreasing (script
\texttt{chantierH\_supnorm.py}). The counter-example (4.23) in \cite{BourgainChang2006},
$\theta=1+p$ (order $1\bmod p$, hence no cancellation), is exactly why Type~A
($\ep\le L$, small order $\bmod\,p$) is \emph{not} covered by
Theorem~\ref{thm:bc}: the Type~A/Type~B split is forced, Type~A being the
García--Voloch regime (Proposition~\ref{prop:gv}) and Type~B the Bourgain--Chang one.

One might hope Theorem~\ref{thm:bc} closes (TB). It does not, for a structural
reason worth recording.

\begin{proposition}[The sup-norm cannot close the conversion]\label{prop:block}
Expanding $E_k^{\mathrm{tot}}$ by additive orthogonality gives, per prime,
$\sum_rN_p(r)^k=(p^2)^{1-k}\sum_{\xi_1+\dots+\xi_k\equiv0}\prod_iS(\xi_i)^2$.
Bounding every nonzero frequency by $|S|<B:=(L+1)^{1-\varepsilon_0}$ leaves each
extra nonzero frequency a factor $B^2=(L+1)^{2-2\varepsilon_0}>1$ \textup{(}for any
$\varepsilon_0<1$\textup{)}: the all-nonzero term is $\gtrsim
(L+1)^{2k-2-2(k-2)\varepsilon_0}L^2/(2^kk!\log L)$, which at $k\sim k^\ast$ is
$\exp\!\big(c(1-\varepsilon_0)\log^2L/\log\log L\big)\gg L^2$. Sup-norm alone is
insufficient regardless of $\varepsilon_0$.
\end{proposition}

\begin{proposition}[Circularity: the missing bound \emph{is} (TB)]\label{prop:circ}
$\sum_{\xi\bmod p^2}|S(\xi)|^{2k}=p^2\,E_k^{\mathrm{add}}$, where
$E_k^{\mathrm{add}}=\#\{(j_i,j_i')_{i\le k}:\sum_i2^{j_i}\equiv\sum_i2^{j_i'}\bmod p^2\}$
is the additive energy of $\{2^j\}$ -- exactly the quantity $E_k^{\mathrm{tot}}$
measures. Thus the moment ($L^{2k}$) bound that would close the conversion of
Proposition~\ref{prop:block} is logically equivalent to (TB) itself. The
sup-norm ($L^\infty$) is available and unconditional; the missing $L^{2k}$ bound
is the conjecture. Hence (TB) is a genuine open estimate, not a statement awaiting
the right reference.
\end{proposition}

% ============================================================================
\section{An unconditional theorem for almost all $n$}\label{sec:almost}
% ============================================================================

This section proves a genuine, unconditional theorem about the same range
$p\in(L,L^2]$, by trading the (currently unavailable) worst-case control of
\S\ref{sec:M} for an average-case control, via a second moment and Markov's
inequality. We state its scope precisely: it controls $p\in(L,L^2]$ only: see
the discussion at the end of this section and \S\ref{sec:open} for what this
does, and does not, give for the full two-powers variant.

\begin{proposition}[Half for free]\label{prop:Qfree}
For every odd $n$, $\#\{p\in(L,L^2]:N_p(n)\ge1\}\le\pi(L^2)=o(L^2)$,
unconditionally (there are simply not enough primes in $(L,L^2]$ for the
mere \emph{count} of contributing primes to reach $L^2$).
\end{proposition}

\begin{definition}
$\displaystyle C(n):=\sum_{p\in(L,L^2]}N_p(n)\bigl(N_p(n)-1\bigr)$ -- the
\emph{coincidence term}: twice the number of triples $(p,P,P')$ with
$P\ne P'$ two pairs in $\{0,\dots,L\}^2$ both summing to $n\bmod p^2$.
\end{definition}

By Proposition~\ref{prop:Qfree}, $\sum_{p\in(L,L^2]}N_p(n)\le\pi(L^2)+C(n)$
for every $n$, so it remains to control $C(n)$.

\begin{proposition}[Second moment]\label{prop:secondmoment}
Let $\overline C$ denote the average of $C(n)$ over residues $n$ modulo
$\prod_{p\in(L,L^2]}p^2$ (equivalently, over a full period in each modulus
separately, by CRT). Then
\[
\overline C=\sum_{p\in(L,L^2]}\frac{\Delta_p}{p^2},\qquad
\Delta_p:=\#\{(P,P'):P\ne P',\ \mathrm{sum}(P)\equiv\mathrm{sum}(P')\!\!\pmod{p^2}\}.
\]
For $p\in(L,L^2]$, $\Delta_p\approx\binom{L+2}{2}^2/p^2$ for a quasi-random
set of $\binom{L+2}{2}$ sums in $\Z/p^2$ (confirmed numerically, no
exceptional clustering observed), giving
\[
\overline C\ \approx\ \frac{(L+2)^4}{4}\sum_{p>L}\frac1{p^4}
\ \approx\ \frac{L^4}4\cdot\frac1{3L^3\log L}\ =\ \frac{L}{12\log L}\ =\ o(L^2).
\]
\end{proposition}

\begin{remark}
The approximation $\Delta_p\approx(\#\mathrm{pairs})^2/p^2$ is the
quasi-random heuristic for collision counts and is \emph{not} proved here in
full rigor for every $p$; it is confirmed numerically
($\overline C\approx0.44,0.43,0.51,0.96,0.52$ for $L=16,20,24,28,32$, matching
the predicted $O(L/\log L)$ growth closely). What \emph{is} fully rigorous,
without this heuristic, is the cruder bound obtained by bounding $\Delta_p$
trivially by $\binom{L+2}{2}^2$ and summing $1/p^2$ over $p>L$ via Mertens:
this already gives $\overline C=O(L^2\cdot1/L)=O(L)$, which still suffices
for what follows (any bound $\overline C=o(L^2)$ does).
\end{remark}

\begin{theorem}[Almost all $n$, on the range $(L,L^2{]}$]\label{thm:almost}
For every $\eta>0$, the density of odd $n\le2^{L+1}$ for which
$\sum_{p\in(L,L^2]}N_p(n)\ge\pi(L^2)+\eta L^2$ is at most
$\overline C/(\eta L^2)=O(1/(L\log L))\to0$. In particular, taking
$\eta=0.4$: for all but a density-$O(1/(L\log L))$ exceptional set of odd
$n$,
\[
\sum_{p\in(L,L^2]}N_p(n)\ <\ \pi(L^2)+0.4\,L^2\ =\ (0.4+o(1))(L+1)^2.
\]
\end{theorem}
\begin{proof}
Markov's inequality applied to $C(n)\ge0$. It remains to justify that the mean
of $C(n)$ over odd $n\le2^{L+1}$ equals $\overline C$ (defined by CRT over the
astronomically larger modulus $\prod_{p\in(L,L^2]}p^2$) up to $o(L^2)$. Each
condition $p^2\mid n-2^l-2^m$ meets the interval $[2^L,2^{L+1})$ in a fixed
residue class, with frequency $\tfrac1{p^2}+O(2^{-L})$; summing the $O(2^{-L})$
errors over the $\le\binom{L+2}{2}$ pairs and the $\le\pi(L^2)$ primes gives a
total discrepancy $\ll L^2\cdot\pi(L^2)\cdot2^{-L}\ll L^6\,2^{-L}=o(1)$ between
the interval mean of $C(n)$ and $\overline C$. (Restricting to odd $n$ only
halves the interval and changes nothing, since $2^l+2^m$ is even and $p^2$ is
odd.) Hence the interval mean is $\overline C+o(1)=o(L^2)$, and Markov applies.
\end{proof}

\subsection{What this gives, and what it does not, for the full variant}

Combining Theorem~\ref{thm:complete}, Proposition~\ref{prop:R1} (closing
$p\le L+1$, unconditionally, for \emph{every} $n$, contributing
$(0.3205+o(1))T$) with Theorem~\ref{thm:almost} (closing $p\in(L,L^2]$, for
\emph{almost all} $n$, contributing $(0.4+o(1))T$) gives, for that same
density-$1-O(1/(L\log L))$ set of $n$:
\[
\sum_{\substack{p\le\sqrt n\\ p\le L+1\ \mathrm{or}\ p\in(L,L^2]}}N_p(n)
\ <\ (0.7205+o(1))\,T,
\]
leaving an explicit margin of $\approx0.28\,T$ before the threshold
$T$ of Proposition~\ref{prop:star}. It remains to control the range
$p\in(L^2,\sqrt{n/2}]$ (wall (B) of \S\ref{sec:open}); for that range, on
average, a \emph{first} moment already suffices, as we now show -- so the
almost-all-$n$ statement is in fact complete.

\subsection{Wall (B) on average: a first-moment bound}\label{sec:Bavg}

\begin{theorem}[Wall (B) on average: first moment at a raised threshold]\label{thm:Bavg}
For $N=2^L$, the mean over $n\in[N,2N)$ of
$B(n):=\sum_{p\in(L^2,\sqrt{n/2}]}N_p(n)$ is $O(1/\log L)$. Consequently, by
Markov, for every $\eta>0$,
\[
\#\bigl\{n\in[N,2N):\ B(n)\ge\eta L^2\bigr\}
\ \le\ N\cdot O\!\left(\frac1{\eta L^2\log L}\right)+o(N).
\]
Two useful thresholds: at $\eta L^2\to1$ (i.e.\ threshold $1$), all but a
density-$O(1/\log L)$ set of $n$ have \emph{no} pair $(l,m)$ with $n-2^l-2^m$
divisible by a square $p^2$, $p\in(L^2,\sqrt{n/2}]$; at any fixed $\eta>0$, all
but a density-$O(1/(L^2\log L))$ set have $B(n)<\eta L^2$. The latter is what the
$0.28\,T$ margin permits, and it is a factor $L^2$ smaller than the former.
\end{theorem}
\begin{proof}
The mean is bounded by summing over pairs and primes:
\[
\sum_{n\in[N,2N)}B(n)\le\sum_{(l,m)}\sum_{p>L^2}
\#\{n\in[N,2N):n\equiv2^l+2^m\!\!\pmod{p^2}\}
\le(L+1)^2\Bigl(N\!\sum_{p>L^2}\tfrac1{p^2}+\pi(\sqrt{N/2})\Bigr),
\]
since each residue class modulo $p^2$ meets $[N,2N)$ in at most $N/p^2+1$
points and the relevant primes are $\le\sqrt{N/2}$. By Mertens,
$\sum_{p>L^2}1/p^2=O(1/(L^2\log L))$, so the first term is
$(L+1)^2\cdot N\cdot O(1/(L^2\log L))=N\cdot O(1/\log L)$; the second is
$(L+1)^2\pi(\sqrt{N/2})=O(L^2\sqrt N/\log N)=o(N)$ since $N\gg L^4$. Dividing by
$N$ gives mean $O(1/\log L)$; Markov's inequality at level $\eta L^2$ gives the
displayed density. (The threshold-$1$ form is the special case, recovering the
union bound.)
\end{proof}

\begin{remark}[Higher moments: a genuine but conditional improvement]
Over a full CRT period $\prod_pp^2$ the counters $N_p(n)$ for distinct $p$ are
exactly independent, so the moments of $B(n)$ and of the range-$(L,L^2]$
collision count $C(n)=\sum_pN_p(N_p-1)$ factor over $p$, and the interval bridge
stays valid for the $2k$-th moment as long as the modulus $(L^2)^{2k}\ll2^L$,
i.e.\ $k\ll L/\log L$. This suggests exceptional density $O(L^{-A})$ for every
fixed $A$. We caution, however, that this is \emph{not} unconditional: the
variance is $\mathrm{Var}(C)=\sum_p\mathrm{Var}(N_p(N_p-1))\le M_{\max}^2\,
\overline C$ where $M_{\max}=\max_pM_p$, and with only the trivial
$M_{\max}\le L+1$ this gives $\mathrm{Var}(C)=O(L^2\overline C)$ and Chebyshev
\emph{no} gain over Markov; a genuine gain (numerically $\mathrm{Var}(C)/
\overline C\approx2.1$--$2.9$ to $L=36$, giving density $O(1/(L^3\log L))$)
requires $M_{\max}=o(L)$, or more precisely an unconditional bound on the
fourth-moment sum $\sum_p p^{-2}\sum_r[c_r(c_r-1)]^2$ -- a higher square-sieve
estimate of the same nature as~(TB). So the moments beyond the first are
controlled by the same open input as (M$''$) itself; only the first moment
(Markov) is unconditional. The wide range (B), where $N_p\in\{0,1\}$ generically,
is closer to Bernoulli and its second moment is more plausibly unconditional, but
we do not pursue it. The first moment already suffices for the corollary.
\end{remark}

This is the structural counterpart of the impossibility noted in
\S\ref{sec:open}: there is simply not enough mass in $\sum_{p>L^2}1/p^2$
(which is $<\sum_{p\ge3}1/p^2=0.2022\ldots$) for the wide range (B) to matter
on average, even though it is the numerically dominant range. Numerically
(script \texttt{checkB.py}), the fraction of $n$ with a nonzero (B)-contribution
is $0\%,3.4\%,5.4\%$ at $L=15,18,22$, average contribution $\le0.056$, in line
with $(L+1)^2\sum_{p>L^2}1/p^2\approx0.13$ (slowly decreasing).

\subsection{The perfect-square range, for almost all $n$}

For the every-$n$ statement the range $(\sqrt{n/2},\sqrt n]$ is closed by
Kalinin~\cite{Kalinin2026} (\S\ref{sec:open}). For the \emph{almost-all}
statement it needs no such input: a first moment suffices, keeping the
unconditional theorem entirely self-contained.

\begin{proposition}[Perfect-square range, almost all $n$]\label{prop:sqrange}
For all but $O\big((L+1)^2\,2^{-L/2}\big)=o(2^L)$ odd $n\in[2^L,2^{L+1})$, no pair
$(l,m)$ has $n-2^l-2^m$ a perfect square; a fortiori the range
$(\sqrt{n/2},\sqrt n]$ contributes $0$.
\end{proposition}
\begin{proof}
For $p>\sqrt{n/2}$ and $p^2\mid k$ with $0<k<n\le2p^2$, the only multiple of
$p^2$ available is $k=p^2$; so any contribution from this range requires
$n-2^l-2^m$ to be a perfect square. Such an $n$ equals $2^l+2^m+s^2$ for some
$l,m\le L$ and $s\ge0$. The number of triples $(l,m,s)$ with
$2^l+2^m+s^2\in[2^L,2^{L+1})$ is
$\sum_{l,m}\#\{s:s^2\in[2^L-2^l-2^m,\,2^{L+1}-2^l-2^m)\}\le(L+1)^2(2^{L/2}+1)$,
since \emph{any} interval of length $X$ holds at most $\sqrt X+1$ squares
(consecutive squares in $[a,a+X)$ satisfy $\sqrt{a+X}-\sqrt a\le\sqrt X$); here
$X=2^L$, and the interval may start near $0$ when $2^l+2^m\approx2^L$, so the
universal bound is needed. Hence the exceptional set has size
$\le(L+1)^2(2^{L/2}+1)$, of density $O((L+1)^2\,2^{-L/2})\to0$.
\end{proof}

\subsection{The complete almost-all-$n$ theorem}

\begin{corollary}[Two-powers variant for almost all $n$, unconditional]\label{cor:almostall}
For all but a density-$O(1/(L\log L))$ set of odd $n\le2^{L+1}$, the two-powers
variant holds.
\end{corollary}
\begin{proof}
Outside the union of the exceptional sets -- never (for $p\le L+1$, by
Theorem~\ref{thm:complete} and Proposition~\ref{prop:R1}, deterministically
$\le(0.3205+o(1))T$); density $O(1/(L\log L))$ (for $(L,L^2]$, by
Theorem~\ref{thm:almost} at $\eta=0.4$, contributing $<(0.4+o(1))T$); density
$O(1/(L^2\log L))$ (for $(L^2,\sqrt{n/2}]$, by Theorem~\ref{thm:Bavg} at
$\eta=0.1$, contributing $<0.1\,T$); density $O((L+1)^2\,2^{-L/2})$ (for
$(\sqrt{n/2},\sqrt n]$, contributing $0$, by Proposition~\ref{prop:sqrange}) -- we
get $\sum_{p\le\sqrt n}N_p(n)<(0.8205+o(1))T<T$, whence a squarefree pair exists by
Proposition~\ref{prop:star}. The union of exceptional sets has density
$O(1/(L\log L))\to0$, dominated by $(L,L^2]$; in particular the unconditional
theorem now uses no external input for any range.
\end{proof}

Thus the only gaps to a \emph{full} (every-$n$) proof are the worst-case
statements (M$''$) and (B) of \S\ref{sec:open}; the almost-all-$n$ statement
is unconditional and complete. Raising the wall-(B) threshold from $1$ to
$\eta L^2$ (which the $0.28\,T$ margin permits) has moved the binding term in
the exceptional density from (B) to the range $(L,L^2]$, at $O(1/(L\log L))$ --
a factor-$L$ improvement over the naive $O(1/\log L)$. A further conditional
improvement to $O(L^{-A})$ (any $A$) is available via higher moments, but rests on
a fourth-moment bound of the same difficulty as (M$''$) (see the preceding
remark); the unconditional density remains $O(1/(L\log L))$.

% ============================================================================
\section{What remains: the open links}\label{sec:open}
% ============================================================================

We now state, as precisely as possible, the two gaps left open by the above,
together with where each brute-force argument fails and by how much.

\subsection{The range $(\sqrt{n/2},\sqrt n]$: closed, by Kalinin (2026)}

If $p\in(\sqrt{n/2},\sqrt n]$ and $N_p(n)\ge1$ with the further coincidence
$2^l=2^m$ (i.e.\ a perfect-square term $n-2^{l+1}$), Kalinin
\cite{Kalinin2026} proves (his Theorem~4, building on a multiplicity bound
$|\{k\ge1:n-2^k=mu^2\}|\le2$ for squarefree odd $m\ge3$, sharp, via an
elementary $2$-adic and mod-$8$ analysis with no equidistribution input) that
this range contributes only $O(L)$ to the relevant count, for every $n$.
This range is therefore \textbf{closed, unconditionally, for every $n$} --
the only range, besides $p\le L+1$, with that status.

\subsection{Open link (M$''$): the multiplicity-distribution tail on $(L,L^2]$}

Precisely, what is missing to upgrade Lemma~\ref{lem:M} from a deterministic
\emph{reduction} to a deterministic \emph{bound} is:
\[
\text{(M$''$)}\qquad
\sum_{p\in(L,L^2]\ \mathrm{non\text{-}Sidon}}\bigl(M_p-2\bigr)\ =\ o(L^2),
\qquad\text{equivalently}\qquad
\sum_{k\ge2}\#\{p\in(L,L^2]:M_p^{\mathrm{unord}}\ge k\}\ =\ o(L^2).
\]
\textbf{Where brute force fails, exactly.} The trivial bound
$M_p-2\le L-1$ together with the trivial bound $\le\pi(L^2)$ on the number of
non-Sidon primes gives $\Theta(L^2)$, not $o(L^2)$ -- it fails by a constant
factor, not by an order of magnitude. \textbf{What the data shows.}
$M_p^{\mathrm{unord}}$ (the maximum number of \emph{unordered}
representations of any residue $n\bmod p^2$ as $\{2^i,2^j\}$) is \emph{not}
bounded by an absolute constant: it grows slowly, empirically
$3,4,4,4,5,6$ for $L=20,30,40,50,60,70$, attained on primes with anomalously
small multiplicative order of $2$ (e.g.\ $p=127$, $\ord_p(2)=7$). We checked,
and explicitly ruled out, the natural attempt to bound $M_p$ via Evertse's
$S$-unit equation theorem: that theorem bounds non-degenerate solutions of
$S$-unit equations over number fields (characteristic $0$); it does not
apply to equations in the finite ring $\Z/p^2$, where $S$-unit-type equations
$x+y=z+w$ over a multiplicative subgroup $H$ typically have
$\sim|H|^2/p^2\cdot|H|$ solutions, unbounded as $|H|\to\infty$ -- the
opposite regime from characteristic $0$. The proportion of non-Sidon primes
in $(L,L^2]$ is itself a constant fraction ($\approx17\%$, stable from
$L=32$ to $L=48$), not a vanishing one -- so (M$''$) is genuinely a
statement about the \emph{decay of the multiplicity distribution} (a
$B_2[g]$-type statement: how many primes have additive multiplicity $\ge k$,
for each $k$), not merely about how many primes fail to be Sidon. \textbf{A cleaner equivalent form.} Since $\#\{$non-Sidon $p\in(L,L^2]\}\le
\pi(L^2)=O(L^2/\log L)$ unconditionally, writing the average excess
$\mathrm{EM}(L):=\bigl(\sum_{\text{non-Sidon}}(M_p-1)\bigr)/\#\{\text{non-Sidon}\}$
gives the equivalence
\[
\text{(M$''$)}\quad\Longleftrightarrow\quad \mathrm{EM}(L)=o(\log L).
\]
Empirically $\mathrm{EM}(L)\approx1.3$ (computed to $L=120$: it lies in
$[1.27,1.33]$, with $\mathrm{EM}(L)/\log L$ \emph{decreasing}, $0.33\to0.27$),
so (M$''$) holds with a growing margin; correspondingly $\sum_pM_p/L^2$
decreases ($0.16\to0.135$). We had conjectured $\mathrm{EM}\to4/\pi=1.273$ from
data to $L=80$; extending to $L=120$ \emph{refutes} a clean closed form
(the value drifts up to $\approx1.30$). What is robust is only
$\mathrm{EM}=O(1)$, indeed $o(\log L)$. \textbf{The shape of the
distribution.} The reason $\mathrm{EM}=O(1)$ holds is that the tail counts
$\#\{p:M_p^{\mathrm{unord}}\ge k\}$ decay \emph{geometrically} in $k$: on the
Type~B primes (\S\ref{sec:classify}) the ratio
$\#\{M_p\ge k{+}1\}/\#\{M_p\ge k\}$ is $\approx0.08$ and stable across
$L=60,\dots,180$ ($n$ up to $\sim10^{54}$), with $\#\{M_p\ge5\}=0$ throughout;
the excess $\sum_k\#\{M_p\ge k\}$ is thus carried almost entirely by the level
$k=3$, and $\#\{M_p\ge2\}\le\pi(L^2)=o(L^2)$ is unconditional. We stress that
this geometric decay is \emph{observed}, not proved inductively, and that it is
the empirical \emph{shape} of the same distribution, not an independent
simplification: proving the decay uniformly in $k$ is exactly proving the
sub-Poisson behaviour below, and does not sidestep it.

\textbf{Why four standard toolkits all fail, for one reason.} We tried, and
ruled out, the four natural routes; each controls an \emph{energy}-scale
quantity, whereas (M$''$) needs the \emph{maximum multiplicity}, and for these
near-Sidon geometric sets the maximum is far below the energy:
\begin{itemize}
\item \emph{Cauchy--Schwarz / additive energy:} $M_p-1\le\sqrt{2\Delta_p}$,
but $\sum_p\sqrt{2\Delta_p}=\Theta(L^2)$ (ratio $\approx0.25$, not $\to0$); and
$\sum_p\Delta_p$ itself \emph{grows} faster than $L^2$.
\item \emph{$S$-unit / Evertse:} inapplicable in the finite ring $\Z/p^2$
(as above; $S$-unit-type equations there have $\sim|H|^3/p^2$ solutions).
\item \emph{Plünnecke--Ruzsa:} the geometric set $G_L=\{2^k\bmod p^2\}$ has
\emph{maximal} doubling $|G_L+G_L|/|G_L|\sim L$, so Plünnecke gives only the
trivial $E\le K^2|G_L|^2\sim L^4/4$. (Yet $E/|G_L|^2\approx2$--$4$: $G_L$ is
``almost Sidon'', energy barely above the floor.)
\item \emph{Exponential sums (Bourgain--Garaev--Konyagin--Shparlinski):} a
pointwise bound on $S(t)=\sum_{k\le L}e_{p^2}(t2^k)$ bounds the energy
$E_p=p^{-2}\sum_t|S(t)|^4$, not $M_p$, and $M_p\ll\sqrt{E_p}$ (measured gap
$M_{127}=9$ vs $\sqrt{2\Delta_{127}}=167$ at $L=120$). Moreover our sums are
incomplete of length $L+1\in[q^{1/4},q^{1/2}]$, $q=p^2$, where even Korobov's
$\sqrt q\log q$ is trivial -- below the range of known unconditional cancellation
for the exponential function mod $p^2$. \emph{A $p$-adic reduction (Lemma~K
applied to the moment, not just the main term) partly clarifies where the
difficulty sits}: splitting $S(t)$ by whether $p\mid t$, the $p\mid t$
characters $t=pb$ give $S(pb)=\sum_{k\le L}e_p(b\,2^k)$ -- a \emph{complete} sum
over $\langle2\rangle\bmod p$ when $\ep\le L$ (Type~A, already unconditional by
Proposition~\ref{prop:ordersum}), an incomplete geometric sum mod~$p$ otherwise
-- whereas the $p\nmid t$ characters give genuinely incomplete sums mod~$p^2$
that Lemma~K does \emph{not} reduce. So the true frontier is not Korobov's
generic bound but \emph{sum-product} cancellation for the geometric progression
itself (Bourgain--Glibichuk--Konyagin: incomplete geometric sums of length $L$
mod~$p$ or~$p^2$, valid for order $\ep\ge p^\delta$); whether such a bound
reaches length $L$ in this range, in mean square, is precisely the open input
(TB), which we have not settled.
\end{itemize}
Equivalently (double counting), $\sum_p\Delta_p=\sum_{\text{nontrivial }(a,b,c,d)}
\#\{p:p^2\mid N\}$ with $N=2^a+2^b-2^c-2^d\ne0$ (the latter is exactly why
$\{2^k\}$ is a Sidon set over $\Z$); here $\#\{p\in(L,L^2]:p^2\mid N\}\le2$ in
range, and the only $\omega_2(N)=2$ cases are the structured squares
$N=(2^k-1)^2$ (e.g.\ $(2^{14}-1)^2=3^2\!\cdot\!43^2\!\cdot\!127^2$). All of this
controls the \emph{column/row} marginals (per-quadruplet, per-prime energy) but
never the \emph{pile} marginal $M_p$ (e.g.\ for $p=127$, $L=40$: row $\Delta_p=329$
collisions fan into piles of max $M_p=3$). So (M$''$) is, precisely, a
\emph{large-deviation / multiplicity-distribution} statement about the geometric
progression $\{2^k\bmod p^2\}$ -- the maximum pile, not the energy -- which lies
beyond all four \emph{energy}-scale toolkits above. (A fifth, maximum-scale
family -- Stepanov / García--Voloch -- does bound the pile directly
(Proposition~\ref{prop:gv}), but on Type~A alone and short of the
target by a power of $L$.) We regard it as true (every empirical margin
grows) but accessible only to a dedicated argument of that kind, which we do not
have. As a structural aside, the high-$M_p$ primes correlate strongly with
\emph{small} multiplicative order $\ord_p(2)$ ($\mathrm{corr}(M_p,1/\ord_p(2))
\approx0.71$ at $L=60$), though large-order primes also occur among the
non-Sidon ones, so this does not by itself reduce the sum.

\textbf{The refined picture (\S\ref{sec:classify}--\ref{sec:typeB}).} The
correlation with $\ord_p(2)$ is in fact the entry point to a clean split of
(M$''$). Writing Type~A $=\{\ep\le L\}$, Type~B $=\{\ep>L\}$: the Type~A excess is
$\le\kappa_p(\lfloor L/\ep\rfloor+1)$ per prime (Proposition~\ref{prop:perclass}),
and the order sum $\sum_p\lfloor L/\ep\rfloor=o(L^2)$ is now
\emph{unconditional} (Proposition~\ref{prop:ordersum}); the only missing Type~A
input is a bound on the number $\kappa_p$ of active progressions (hypothesis
(TA)), for which a \emph{fifth}, maximum-scale toolkit -- the Stepanov /
García--Voloch method for additive representations in a multiplicative subgroup --
gives the unconditional $\kappa_p\ll\ep^{2/3}$ and hence
$\sum_{\text{Type A}}M_p=O(L^{8/3}/\log L)$, short of $o(L^2)$ by $L^{2/3}$
(Proposition~\ref{prop:gv}). The Type~B excess is governed by
the factorial moments $E_k^{\mathrm{tot}}$ (Lemma~\ref{lem:markov}), and the sharp
condition $E_k^{\mathrm{tot}}\le C^kE_k^{\mathrm{null}}$ (hypothesis (TB),
Theorem~\ref{thm:conditional}) reduces it to a Heath-Brown square-sieve estimate,
i.e.\ a \emph{mean-square} bound on the same incomplete exponential sums that
appear in sup-norm in the fourth toolkit above. Thus (M$''$) is not monolithic:
its order-scale part is unconditional, its Type~A part is a rigidity statement
(now with the unconditional partial bound $\kappa_p\ll\ep^{2/3}$ above), and its
Type~B part is the one genuine analytic input. It is worth stressing how
\emph{singular} this lock is: even the unconditional almost-all density of
\S\ref{sec:almost} cannot be improved past the first moment without it. The
variance of the collision count $C(n)=\sum_pN_p(N_p-1)$ is
$\sum_p\mathrm{Var}(N_p(N_p-1))$, whose per-prime terms are the sums
$\sum_rc_r^{\,4}/p^2$ -- i.e.\ the very moments $E_4^{\mathrm{tot}}$; likewise a
truncation $C^{\le K}+C^{>K}$ pushes the tail onto $\sum_p\sum_r\binom{c_r}{K}$,
again $E_K^{\mathrm{tot}}$. So Chebyshev (numerically a factor-$L^2$ gain to
$O(1/(L^3\log L))$, $\mathrm{Var}(C)/\overline C\approx2.1$--$2.9$ to $L=36$) and
every higher-moment sharpening are conditional on the \emph{same} square-sieve
input as (M$''$) itself. There is one underlying lock, and it is the
multiplicity distribution of $\{2^l+2^m\bmod p^2\}$.

\subsection{Open link (B): the equidistribution wall on $(L^2,\sqrt{n/2}]$ --
worst case only}

\emph{On average} this range is already handled (Theorem~\ref{thm:Bavg}: its
first moment is $O(1/\log L)$, so it contributes nothing for almost all $n$).
What remains genuinely open is the \emph{uniform-in-$n$} (worst-case) statement:
\[
\text{(B)}\qquad
\sum_{p\in(L^2,\sqrt{n/2}]}N_p(n)\ =\ o(L^2)\quad\text{uniformly in }n.
\]
\textbf{Where brute force fails, exactly.} For $p$ in this range,
$N_p(n)\in\{0,1\}$ generically (no two pairs can plausibly collide once $p$
is this large relative to the number of pairs), so the sum is, up to a
constant, the number of pairs $(l,m)$ for which $n-2^l-2^m$ has \emph{some}
square prime factor $p^2$ with $L^2<p^2\le n/2$. Each $k=n-2^l-2^m\le n$ has
at most $O(L/\log L)$ such prime factors (since $(L^2)^{L/\log L}\sim n$), so
the trivial bound is $(L+1)^2\cdot O(L/\log L)=O(L^3/\log L)$ -- the
\emph{cube}, not the square, of $L$: this range is, in the crudest sense, the
\emph{numerically dominant} one (most primes up to $\sqrt n$ lie here), even
though the data shows its actual contribution is small. \textbf{What would
close it.} A squarefree-density / equidistribution estimate for the
structured sequence $\{n-2^l-2^m\}$ in this range -- in spirit close to (but
not identical to) the kind of input behind the classical ``almost all $n$''
result for the one-power variant of \#11 itself. The covering-congruence
machinery that produces the large exceptional sets for the \emph{prime}-kernel
problems \#9 and \#16 (Crocker \cite{Crocker1971}, Pan \cite{Pan2011}, Chen
\cite{Chen2024}) is \emph{provably unavailable} here: a covering system needs
its moduli to satisfy $\sum1/m_i\ge1$, whereas square moduli $p^2$ ($p\ge3$)
give only $\sum_{p\ge3}1/p^2=0.2022\ldots<1$, so one cannot cover $\Z$ by
congruences mod $p^2$. This is the quantitative form of the Crocker decoupling
of \S\ref{sec:intro-decoupling} -- and the same smallness ($\sum_{p>L^2}1/p^2$
tiny) is exactly what makes the average bound Theorem~\ref{thm:Bavg} work. The
worst case (B) would, by contrast, require a genuine equidistribution/large-sieve
estimate for the structured sequence $\{n-2^l-2^m\}$; we have not found one.

We note that (WB) is of the \emph{same nature} as, and in fact weaker than, the
Type~B hypothesis (TB) of \S\ref{sec:typeB}: both are factorial-moment (grand
sieve) bounds on the multiplicity distribution of $\{2^l+2^m\bmod p^2\}$, (WB)
over the larger primes $p>L^2$ (where collisions are rarer, so the maximum
$\max_nB(n)$ is smaller). The naive hope that $\max_nB(n)=O(1)$ is \emph{false}:
$p\le\sqrt{n/2}\le2^{L/2}$ gives $p^2\le2^L$, while a difference of two pair-sums
has absolute value up to $4\cdot2^L>p^2$, so the ``at most one pair per fibre''
argument does not apply in this range; numerically $\max_nB(n)$ tracks the
Poisson maximum $\approx0.69\,L/\log L$ (not $O(1)$), which is nonetheless
$o(L^2)$, so (WB) holds with a wide margin. Thus (TA), (TB) and (WB) are three
scales of a single underlying question -- the multiplicity distribution of
$\{2^l+2^m\bmod p^2\}$ -- rather than three independent walls.

\subsection{Summary table}

\begin{center}
\begin{tabular}{lll}
\toprule
Range of $p$ & Status (every $n$) & Method\\
\midrule
$p\le L+1$ & \textbf{closed} & Lemma K + Prop.~\ref{prop:R1} (elementary)\\
$(L,L^2]$, Type A ($\ep\le L$) & \emph{open}, gap $L^{2/3}$ & $O(L^{8/3}/\log L)$, García--Voloch (Prop.~\ref{prop:gv})\\
$(L,L^2]$, Type B ($\ep>L$) & \emph{open} ((TB); sup-norm route closed by circularity) & (TB) $=$ additive energy (Prop.~\ref{prop:circ})\\
$(\sqrt{n/2},\sqrt n]$ & \textbf{closed} & Kalinin~\cite{Kalinin2026}; a.a.\ $n$ self-contained (Prop.~\ref{prop:sqrange})\\
$(L^2,\sqrt{n/2}]$ & \emph{open}, worst case only & wall (B); \#11-calibre (avg.\ done, Thm~\ref{thm:Bavg})\\
\bottomrule
\end{tabular}\\[2pt]
(All four ranges are \textbf{closed unconditionally for almost all $n$},
Corollary~\ref{cor:almostall}.)
\end{center}

% ============================================================================
\section{The function-field model: where the difficulty really lives}\label{sec:ff}
% ============================================================================

To locate the obstruction precisely we run the whole architecture over
$\F_q[t]$. Replace $\Z$ by $\F_q[t]$, the base $2$ by the variable $t$,
primes $p$ by monic irreducibles $P$, $|p|$ by $|P|=q^{\deg P}$, and
``squarefree'' by squarefree in $\F_q[t]$ (density $1-1/q$). The variant reads:
every $n(t)$ of large degree is $k(t)+t^l+t^m$ with $k$ squarefree. The sieve is
identical, with $N_P(n)=\#\{(l,m):P^2\mid n-t^l-t^m\}$ and threshold
$T=(L+1)^2$. Lemma~K transfers verbatim (the kernel $K_P=1+P\,\F_q[t]/P^2$ of
$(\F_q[t]/P^2)^\times\to(\F_q[t]/P)^\times$ absorbs the character sum), as does
the involution $1+t^{\,e-\delta}\equiv t^{-\delta}(1+t^\delta)$; and
$\sum_P|P|^{-2}=\sum_d(\#\{\deg P=d\})q^{-2d}<1$ (for $q=2$, $\approx0.60$), so
covering by congruences $\bmod\,P^2$ is impossible -- the Crocker decoupling
holds here too.

The decisive difference is at (TB).

\begin{proposition}[{(TB) over $\F_q[t]$, by degree}]\label{prop:ff}
Let $q>2k$. For an irreducible $P$ with $\deg P>L/2$, the sums
$\{t^l+t^m\bmod P^2\}$ have no non-trivial coincidence, so $N_P(r)\le1$ for the
relevant residues and the $k$-th factorial moment vanishes. Consequently the
$\F_q[t]$-analogue of (TB) holds unconditionally.
\end{proposition}
\begin{proof}
A coincidence $\sum_it^{a_i}\equiv\sum_it^{b_i}\pmod{P^2}$ means $P^2$ divides the
polynomial $D=\sum_it^{a_i}-\sum_it^{b_i}$, of degree $\le L$ and with integer
coefficients reduced mod $q$. As $q>2k$, $D$ vanishes in $\F_q[t]$ iff the
exponent multisets coincide. If they do not, $D\ne0$ has $\deg D\le L<2\deg P=\deg
P^2$, so $P^2\nmid D$. Hence the only coincidences are trivial. For $\deg P\le
L/2$ (finitely many $P$, $|P|\le q^{L/2}$) the coincidences are counted by degree,
without any exponential-sum input.
\end{proof}

So over $\F_q[t]$ the analogue of (TB) is a \emph{theorem}, but proved by the
linear independence of the monomials $t^j$ and a degree count -- not by Weil, and
not by any sup-norm bound. Indeed the moment/sup-norm obstruction of
\S\ref{sec:bc} persists here: Weil gives square-root cancellation
$\varepsilon=\tfrac12$, still $B^2=(L+1)>1$, so the conversion of
Proposition~\ref{prop:block} blocks identically; the closure comes from
elsewhere. And the mechanism does \emph{not} transpose to $\Z$: there is no
degree bound preventing $p^2\mid2^a+2^b-2^c-2^d\ne0$. The obstruction over $\Z$ is
thus \emph{archimedean} -- powers of $2$ carry (e.g.\ $2^a+2^a=2^{a+1}$) while
monomials $t^a$ do not, so $\{2^j\}$ is only \emph{approximately} Sidon and a
nonzero four-term combination can still be squarefull -- rather than a deficit of
cancellation. The function-field model therefore supplies no characteristic-zero
method; it certifies that (TB) over $\Z$ is genuinely harder, of square-sieve /
squarefull-value type.

% ============================================================================
\section{Conclusion}
% ============================================================================

The two-powers squarefree variant of Erd\H{o}s \#11 is, as far as we have
been able to push it, reduced to exactly two clean, isolated, named
statements -- (M$''$), a multiplicity-distribution decay on a bounded range
of primes just above $L=\log_2n$, and (B), a squarefree-density statement on
the much wider range of primes between $L^2$ and $\sqrt n$ -- after an
elementary closure (Lemma~K and Proposition~R1) that handles, exactly and for
every $n$, every prime up to $L+1$, and a deterministic dissolution of the
worst-case obstruction (Lemma~M) on the dominant range $(L,L^2]$. An
unconditional theorem (\S\ref{sec:almost}--\ref{sec:Bavg}) further proves the
variant for \emph{almost all} $n$, on \emph{all} prime ranges
(Corollary~\ref{cor:almostall}).

Within (M$''$) the split of \S\ref{sec:classify}--\ref{sec:typeB} is now sharp.
The order-scale part is unconditional (Proposition~\ref{prop:ordersum}). The
Type~A part carries a genuine \emph{unconditional} bound $\kappa_p\ll\ep^{2/3}$
via García--Voloch (Proposition~\ref{prop:gv}), giving $\sum_{\text{Type A}}M_p=
O(L^{8/3}/\log L)$ -- short of the target $o(L^2)$ by exactly $L^{2/3}$, so (TA)
survives only as the sharp question of pushing $2/3\to o(1)$ using the involution
rigidity. The Type~B part (TB) is where the problem's true nature surfaces. The
relevant \emph{additive} exponential sums do satisfy an unconditional sup-norm
bound $|S(\xi)|<(L+1)^{1-\varepsilon}$ (Bourgain--Chang, Theorem~\ref{thm:bc},
confirmed numerically), yet this cannot close (TB): the moment ($L^{2k}$) bound
one would need is, by Proposition~\ref{prop:circ}, \emph{identical} to the
additive energy $E_k^{\mathrm{add}}$, i.e.\ to (TB) itself. (TB) is therefore
closed by \emph{circularity}, not by any absent reference -- it is a genuine open
estimate. The function-field model (\S\ref{sec:ff}) pins down why: the analogue
of (TB) over $\F_q[t]$ is a theorem, but proved by the linear independence of the
monomials $t^j$ and a degree count, a mechanism that does \emph{not} transpose
because over $\Z$ nothing bounds the squarefull part of a nonzero four-term
combination $2^a+2^b-2^c-2^d$. The obstruction is \emph{archimedean} -- powers of
$2$ carry while monomials do not -- not a deficit of cancellation. In this precise
sense (M$''$) at $k=2$ is a squarefree-value statement for the sparse family
$1+2^\beta-2^\gamma-2^\delta$, a genuine cousin of \#11 itself; empirically the
few quadruples that a single $p^2$ divides twice are \emph{all} structured
(Mersenne squares $N=(2^k-1)^2$ and relatives; \S\ref{sec:open}, script data),
suggesting an elementary classification for $k=2$ that we state as a conjecture,
not a theorem. As for (B): its average form is settled (first moment); its
worst case remains, and we judge it to be of the same calibre as \#11 in its
original one-power form. The
\emph{elementary} lemmas underpinning the unconditional results
(Proposition~\ref{prop:ordersum}, the order-sum bound; Lemma~\ref{lem:involution},
the involution identity; the order lower bound behind
Proposition~\ref{prop:convergence}; the coprimality product bound behind
Proposition~\ref{prop:ordersum}; and the degree obstruction behind the
function-field Proposition~\ref{prop:ff}) have been \emph{formally verified in
Lean~4} against Mathlib (file \texttt{Erdos11\_verified.lean}, five lemmas,
\texttt{0 sorry}, axiom audit \texttt{[propext, Classical.choice, Quot.sound]});
the analytic and
conditional statements (the factorial-moment bound (TB), the every-$n$ theorem)
are \emph{not} formalised, matching their status as conjectural. Only the target
statement \texttt{erdos\_11.variants.two\_pow\_two} otherwise exists in Lean form.
We hope the precision of the named gaps is useful to whoever picks this up next.

% ============================================================================
\section*{Acknowledgements}
This note was produced in a human-directed workflow with AI assistants in
distinct executor and validator roles: proposing tools and computations,
locating and cross-checking references, and deriving and refuting bounds, with a
standing rule that no theorem enters the paper without a written proof or a
formal (Lean) verification, and every numerical assertion points to a named
script (see the repository \texttt{README}). The mathematical judgements,
and the responsibility for errors, are the author's.

% ============================================================================
\begin{thebibliography}{9}

\bibitem{erdosproblems11}
T.~F.~Bloom, \emph{Erd\H{o}s Problem \#11},
\url{https://www.erdosproblems.com/11}.

\bibitem{Kalinin2026}
D.~Kalinin, \emph{A multiplicity bound for a fixed squarefree kernel}, preprint,
29 March 2026. We rely on this paper for the exact statements of, and exact
references to, Granville--Soundararajan, Hercher, Crocker and
Platt--Trudgian on the one-power variant of \#11, and for Theorem~4 used in
\S\ref{sec:open}.

\bibitem{Crocker1971}
R.~Crocker, \emph{On the sum of a prime and two powers of two}, Pacific J.
Math. \textbf{36} (1971), 103--107. (Title and pages confirmed via the
bibliography of \cite{Chen2024} and \cite{Pan2011}.)

\bibitem{Chen2024}
Y.-G.~Chen, \emph{A conjecture of Erd\H{o}s on $p+2^k$}, arXiv:2312.04120v3
(2024) -- the resolution of Erd\H{o}s Problem~\#16 (the \emph{one}-power
exceptional-set structure), used here only to fix the exact references and the
covering-congruence mechanism discussed in \S\ref{sec:open}.

\bibitem{Pan2011}
H.~Pan, \emph{On the integers not of the form $p+2^a+2^b$}, Acta Arith.
\textbf{148} (2011), 55--61 (arXiv:0905.3809) -- the \emph{prime}-kernel
two-powers problem (Erd\H{o}s \#9); its lower bound on non-representable
integers is obtained by a covering system modulo primes, not a sieve.

\bibitem{HeathBrown1984}
D.~R.~Heath-Brown, \emph{The square sieve and consecutive square-free numbers},
Math. Ann. \textbf{266} (1984), 251--259.

\bibitem{GarciaVoloch1988}
A.~García, J.~F.~Voloch, \emph{Fermat curves over finite fields}, J. Number
Theory \textbf{30} (1988), 345--356 (primary source of the multiplicity bound
$r_p(b)\le4|G|^{2/3}$).

\bibitem{Konyagin-notes}
S.~V.~Konyagin, \emph{Estimates for character sums and the additive structure of
multiplicative subgroups} (lecture notes; recordings and notes available via
\url{https://www.mathtube.org/}); and S.~V.~Konyagin, I.~E.~Shparlinski,
\emph{Character sums with exponential functions and their applications}, Cambridge
Univ. Press, 1999. We take the exact form used from here (Konyagin's Theorem~2.1):
$r_p(b)\le4|G|^{2/3}$ for $b\in\Z_p^\times$ and $|G|<(p-1)/((p-1)^{1/4}+1)$; the
Stepanov-method proof (Heath-Brown) and the energy bound $T_2(G)\ll|G|^{5/2}$ for
$|G|\le p^{2/3}$ (Heath-Brown--Konyagin) are stated there. The exact size and
non-degeneracy conditions used in Proposition~\ref{prop:gv} are those of this
secondary source, verified numerically herein.

\bibitem{BourgainChang2006}
J.~Bourgain, M.-C.~Chang, \emph{Exponential sum estimates over subgroups and
almost subgroups of $\Z_Q^*$, where $Q$ is composite with few prime factors},
Geom. Funct. Anal. \textbf{16} (2006), 327--366. Corollary~4.5 gives the additive
sup-norm of Theorem~\ref{thm:bc}; formula (4.23) ($\theta=1+p$) is the degeneracy
that forces the Type~A/B split.

\bibitem{Rudin1960}
W.~Rudin, \emph{Trigonometric series with gaps}, J. Math. Mech. \textbf{9}
(1960), 203--227.

\bibitem{erdos11data}
Numerical evidence referenced throughout (scripts \texttt{orders.py},
\texttt{lemmaE2.py}, \texttt{verify\_structure.py}, \texttt{verify\_residual.py},
and further task-specific computations for $M_p$, $C(n)$, and the
non-Sidon-prime counts, all in this repository).

\end{thebibliography}

\end{document}
